\documentclass[12pt,a4paper]{article}

\usepackage[utf8]{inputenc}
\usepackage[T2A]{fontenc}
\usepackage[russian]{babel}
\usepackage{geometry}
\usepackage{graphicx}
\usepackage{hyperref}
\usepackage{longtable}
\usepackage{array}
\usepackage{subcaption}
\geometry{margin=2.5cm}

\title{Chronos --- платформа в стиле Duolingo для изучения истории}
\author{Инютин Святослав, Надточий Ярослав, Карпов Максим}
\date{2026}

\begin{document}

\maketitle

\section{Описание проекта}

Chronos --- образовательная платформа для изучения истории через короткие
интерактивные уроки. Проект ориентирован на формат регулярного обучения:
пользователь проходит небольшие сессии, отвечает на задания разных типов,
получает прогресс по темам и рекомендации, что изучать дальше.

В текущей версии основной учебный материал посвящён истории России начала
XX века: правлению Николая II, реформам С. Ю. Витте, революции 1905--1907 годов
и рабочему вопросу. В дальнейшем планируется расширять объём тем и набор
форматов заданий.

Основная цель проекта --- сделать изучение истории более интерактивным и
вовлекающим за счёт игровых механик, системы прогресса, коротких обучающих
сессий и повторения материала.

\section{Проблема и актуальность}

История часто воспринимается как набор дат, имён и событий, которые нужно
запомнить перед контрольной или экзаменом. Такой подход плохо поддерживает
регулярное повторение и не показывает связи между событиями.

Chronos решает эту проблему за счёт:

\begin{itemize}
    \item дробления материала на короткие темы и задания;
    \item интерактивной проверки знаний после каждого блока;
    \item сохранения прогресса пользователя;
    \item рекомендаций по следующему шагу обучения;
    \item мотивации через XP, streak и достижения.
\end{itemize}

\section{Целевая аудитория}

Проект может быть полезен школьникам, студентам и всем пользователям, которые
хотят изучать историю небольшими регулярными сессиями. Особенно хорошо такой
формат подходит для подготовки к проверочным работам, экзаменам и повторения
сложных тем.

\section{Основной функционал}

\begin{itemize}
    \item регистрация и авторизация пользователей;
    \item просмотр опубликованных исторических курсов;
    \item прохождение коротких learning-сессий;
    \item задания разных типов: теория, один ответ, правда/ложь, заполнить пробел,
          сопоставление и другие;
    \item проверка ответов на серверной стороне;
    \item сохранение попыток пользователя;
    \item прогресс по курсу, разделам, юнитам и навыкам;
    \item XP, счётчик дней подряд и достижения;
    \item рекомендации следующего материала;
    \item административная панель для работы с контентом;
    \item seeder для загрузки стартового курса.
\end{itemize}

% \section{Пользовательский сценарий}

% Основной сценарий работы приложения выглядит следующим образом:

% \begin{enumerate}
%     \item пользователь регистрируется или входит в аккаунт;
%     \item открывает экран обучения и выбирает тему;
%     \item приложение запускает learning-сессию по выбранному skill;
%     \item backend подбирает задания и сохраняет сессию;
%     \item пользователь отвечает на задания;
%     \item learning-модуль проверяет ответы и сохраняет попытки;
%     \item после завершения сессии обновляется progress;
%     \item gamification начисляет XP и обновляет streak;
%     \item recommendation предлагает следующий шаг обучения.
% \end{enumerate}

\section{Архитектура проекта}

Backend проекта построен как модульный монолит. Каждый крупный домен вынесен в
отдельный модуль, а взаимодействие между модулями происходит через интерфейсы.
Это позволяет развивать функциональность независимо и не смешивать бизнес-логику
разных частей приложения.

\subsection*{Основные backend-модули}

\begin{itemize}
    \item \texttt{users} --- пользователи, роли, авторизация и refresh-сессии;
    \item \texttt{content} --- курсы, разделы, юниты, skills и challenges;
    \item \texttt{learning} --- движок прохождения урока и проверки ответов;
    \item \texttt{progress} --- состояние пользователя на карте курса;
    \item \texttt{gamification} --- XP, streak и достижения;
    \item \texttt{recommendation} --- подбор следующего skill и заданий;
    \item \texttt{content/seeder} --- загрузка стартового курса из JSON.
\end{itemize}

\subsection*{Learning-сессия}

Модуль \texttt{learning} реализован как движок урока. Он отвечает за запуск
сессии, выдачу текущего задания, проверку ответа, сохранение попытки и завершение
сессии. При завершении learning передаёт агрегированный результат в
\texttt{progress} и \texttt{gamification}.

\subsection*{Progress}

Модуль \texttt{progress} не знает содержимое заданий и правильные ответы. Он
получает только итог сессии: пользователя, skill, количество правильных ответов,
общее количество ответов и время завершения. На основе этих данных он обновляет
mastery, уровень skill и открывает следующий материал.

\subsection*{Gamification}

Модуль \texttt{gamification} отвечает за мотивационные механики: начисление XP,
обновление streak, уровень пользователя и достижения. Повторная обработка одной
и той же learning-сессии не начисляет XP повторно.

\subsection*{Recommendation}

Модуль \texttt{recommendation} не изменяет состояние пользователя. Он читает
content и progress и определяет, что лучше дать пользователю дальше: продолжить
текущий skill, повторить слабую тему или начать новый skill.

\section{Frontend}

Frontend реализован на Flutter. Приложение содержит экраны авторизации,
главную страницу, карту обучения, страницу урока, квиз, профиль пользователя и
административную панель.

Frontend взаимодействует с backend через REST API. Для интеграции выделены
отдельные API-клиенты:

\begin{itemize}
    \item \texttt{AuthApi};
    \item \texttt{ContentApi};
    \item \texttt{LearningApi};
    \item \texttt{GamificationApi};
    \item \texttt{RecommendationApi}.
\end{itemize}

\section{Технологии}

\begin{itemize}
    \item Go;
    \item PostgreSQL;
    \item REST API;
    \item Flutter;
    \item Docker;
    \item Goose migrations.
\end{itemize}

% \section{База данных и миграции}

% Для хранения данных используется PostgreSQL. Схема базы развивается через
% миграции. Для основных сущностей проекта используется подход ``одна сущность ---
% одна миграция'', что упрощает чтение истории изменений базы данных.

% В базе хранятся пользователи, учебный контент, learning-сессии, очередь заданий,
% попытки ответов, прогресс пользователя, XP-транзакции, streak и достижения.

\section{Стартовый контент}

Для наполнения приложения реализован seeder. Он читает JSON-файл со структурой
курса и создаёт:

\begin{itemize}
    \item курс;
    \item разделы;
    \item юниты;
    \item skills;
    \item challenges.
\end{itemize}

Seeder безопасен для повторного запуска: если курс уже существует, он завершает
работу без изменения существующего контента.

\section{Примеры интерфейса}

\begin{figure}[h]
    \centering
    \begin{subfigure}{0.31\textwidth}
        \centering
        \includegraphics[width=\textwidth]{Screenshot_20260515_234605.png}
        \caption{Страница урока}
    \end{subfigure}
    \hfill
    \begin{subfigure}{0.31\textwidth}
        \centering
        \includegraphics[width=\textwidth]{Screenshot_20260515_234621.png}
        \caption{Материалы урока}
    \end{subfigure}
    \hfill
    \begin{subfigure}{0.31\textwidth}
        \centering
        \includegraphics[width=\textwidth]{Screenshot_20260515_234652.png}
        \caption{Проверка ответа}
    \end{subfigure}
    \caption{Прохождение темы и работа с заданиями}
\end{figure}

\begin{figure}[h]
    \centering
    \begin{subfigure}{0.31\textwidth}
        \centering
        \includegraphics[width=\textwidth]{Screenshot_20260515_234709.png}
        \caption{Задание true/false}
    \end{subfigure}
    \hfill
    \begin{subfigure}{0.31\textwidth}
        \centering
        \includegraphics[width=\textwidth]{Screenshot_20260515_234807.png}
        \caption{Лента событий}
    \end{subfigure}
    \hfill
    \begin{subfigure}{0.31\textwidth}
        \centering
        \includegraphics[width=\textwidth]{Screenshot_20260515_234504.png}
        \caption{Профиль}
    \end{subfigure}
    \caption{Дополнительные экраны приложения}
\end{figure}

\section{Текущий статус}

На текущем этапе реализованы основные модули backend, базовая интеграция с
frontend, загрузка стартового курса и прохождение учебных сессий. Приложение уже
поддерживает полный цикл: пользователь выбирает тему, проходит задания, получает
прогресс и награды.

Дальнейшее развитие проекта может включать:

\begin{itemize}
    \item расширение набора исторических курсов;
    \item улучшение алгоритма рекомендаций;
    \item более сложную систему достижений;
    \item визуальную карту курса;
    \item аналитику ошибок пользователя;
    \item поддержку повторения материала.
\end{itemize}

\section{Ссылка на проект}

% TODO: вставить реальную ссылку на репозпиторий.
GitHub: \url{https://github.com/SIniutin/history-app-backend}

% TODO: вставить ссылку после записи демонстрации.
Видео-демонстрация: \url{https://vk.com/video317694680_456239890}.

\section{Участники проекта}

\begin{longtable}{|p{5cm}|p{8cm}|}
\hline
ФИО & Роль в проекте \\
\hline
Инютин Святослав & backend, модуль learning, архитектура \\
\hline
Надточий Ярослав & frontend, UI/UX \\
\hline
Карпов Максим & генерация контента, исторические материалы \\
\hline
\end{longtable}

\end{document}
